% from main.tex
%%%%%%%%%%  scan idx 272  |  folio 259  %%%%%%%%%%
parfait rel. à \(S\). Mais d'après le théorème de dualité ([H] III 6.7), on a
\[
 f_*\RHom(E,f^!F)\xrightarrow{\sim}\RHom(f_*(E),F).
\]
Comme \(f_*(E)\) est parfait, on a, d'après (I 7.7),
\[
 \RHom(f_*(E),F)\xrightarrow{\sim}(f_*(E))^{\vee}\otimes^{\mathbf L}F,
\]
et l'on gagne grâce à (4.5) (où \(X=Y\) et \(f=\Id_X\)).

\textbf{Corollaire 4.9.1.} Sous les hypothèses de (4.8), si \(f\) est parfait, \(f^!\) envoie \(D(Y)_{S\text{-}\parf}\) dans \(D(X)_{S\text{-}\parf}\).

\emph{Preuve.} Faire \(E=\mathcal O_X\).

\textbf{Corollaire 4.9.2.} Soient \(S\) un schéma noethérien, \(f:X\to S\) un morphisme compactifiable. Posons \(f^!(\mathcal O_S)=K_X\), \(\RHom(\,\cdot\,,K_X)=D_X\). Pour \(E\in\Ob D(f)_{\parf}\), on a \(D_X(E)\in\Ob D(f)_{\parf}\), et la flèche canonique (I 7.2)
\[
 E\longrightarrow D_XD_X(E)
\]
est un isomorphisme.

\emph{Preuve.} Pour la première assertion, appliquer (4.8) avec \(Y=S\), \(F=\mathcal O_S\). Pour la seconde, on se ramène facilement au cas où \(f\) est une immersion fermée. Il suffit alors de prouver que la flèche déduite de la flèche canonique par application de \(f_*\) est un isomorphisme. D'après la dualité globale ([H] III 6.7), on a
\[
 f_*D_X\xrightarrow{\sim}D_Sf_*.
\]
On est donc ramené à montrer que la flèche canonique
\[
 f_*(E)\longrightarrow D_SD_Sf_*(E)
\]
est un isomorphisme pour \(E\) parfait rel. à \(f\), i.e. \(f_*(E)\) parfait, ce qui résulte de la bidualité ordinaire pour les complexes parfaits (I 7.2).

\section*{5. Applications : théorèmes d'échange et de semi-continuité\footnote{Les résultats de ce numéro ne seront pas utilisés dans la suite du Séminaire.}}

\textbf{Lemme 5.1.} Soit \((X,\mathcal O_X)\) un topos annelé (commutatif), notons \(\Mod(X)\) la catégorie des \(\mathcal O_X\)-Modules, \(D(X)=D(\Mod(X))\). Soient \(E\in\Ob D^-(X)\), et \(p\in\mathbb Z\). Considérons les conditions suivantes :
\begin{enumerate}[label=\alph*)]
\item Le
